\documentclass[11pt,a4paper]{article}
\usepackage[margin=2.2cm]{geometry}
\usepackage{fontspec}
\setmainfont{Palatino}
\setmonofont{Menlo}[Scale=0.85]
\usepackage{xcolor,graphicx,enumitem,titlesec,fancyvrb}
\usepackage[most]{tcolorbox}
\usepackage{xurl}
\usepackage[hidelinks]{hyperref}
\newcommand{\sharedlink}{https://8396-193-136-205-165.ngrok-free.app/collab?token=op2.eyJ2IjoyLCJwaWQiOiI5YTRmOTAzYi0xMjk3LTQ3NTYtODgzZi1lY2M5MGYzZTIyNWIiLCJpaWQiOiJYQ2FVMUdxcmw1OGEifQ.vHHb16x_t6_XlUIMaXzCbnULP5xLROo2e8GzB0ZjMzc}
\definecolor{navy}{HTML}{14213D}
\definecolor{gold}{HTML}{B8892B}
\definecolor{paper}{HTML}{F7F4EC}
\hypersetup{colorlinks,urlcolor=gold!80!black,linkcolor=navy}
\titleformat{\section}{\Large\bfseries\color{navy}}{\color{gold}\thesection}{0.6em}{}
\titlespacing*{\section}{0pt}{1.4em}{0.5em}
\setlength{\parindent}{0pt}\setlength{\parskip}{0.5em}
\newtcolorbox{term}{colback=navy,colframe=navy,coltext=white,arc=2pt,
  left=8pt,right=8pt,top=5pt,bottom=5pt,fontupper=\ttfamily\small}
\newtcolorbox{note}[1][Note]{colback=paper,colframe=gold,arc=2pt,boxrule=0.6pt,
  title={#1},coltitle=navy,colbacktitle=paper,fonttitle=\bfseries,
  left=8pt,right=8pt}
\pagestyle{empty}
\begin{document}

\begin{center}
\includegraphics[width=2.3cm]{numenor-logo.png}\\[0.4em]
{\Huge\bfseries\color{navy} Numenor}\\[0.2em]
{\large\color{gold!80!black} Quick install on a Mac}
\end{center}

\vspace{0.3em}
Numenor is a local \LaTeX{} workspace: project files, a code editor and a live PDF preview side by side,
with optional AI help. It runs on your own Mac and opens in your web browser; your projects stay on your computer.
This guide installs \emph{only} Numenor (you do not need Maziar). It takes about 10 minutes plus the TeX download.


\begin{tcolorbox}[colback=gold!12,colframe=gold,arc=3pt,boxrule=1pt,left=10pt,right=10pt,top=6pt,bottom=6pt]
{\large\bfseries\color{navy} No install needed to open the shared project}\par\smallskip
You can open the shared \LaTeX{} project \textbf{EDGE-AI-Overleaf} right now in any web browser --- on a Mac, a PC or a tablet ---
\emph{without installing Numenor}. Click the button, or copy the link into your browser:\par\medskip
\begin{center}
\href{\sharedlink}{\colorbox{navy}{\color{white}\bfseries\ \ Open the EDGE-AI-Overleaf project\ \ }}
\end{center}
{\footnotesize\ttfamily\color{navy!80}\expandafter\url\expandafter{\sharedlink}}\par\smallskip
{\small The project lives on its owner's computer, so the link works while that computer is on and sharing.
If it does not open, ask the owner to send the link again.
Install Numenor (below) only if you want your own projects on your own Mac.}
\end{tcolorbox}

\begin{note}[What you need]
macOS 12 or newer, an internet connection, and the \textbf{Terminal} app
(press \texttt{Cmd + Space}, type \textit{Terminal}, press Return). Copy each dark box into Terminal and press Return.
\end{note}

\section{Install Homebrew (skip if you have it)}
Homebrew installs the other tools with one command each. Check first: if \texttt{brew --version} prints a version, skip this step.
\begin{term}
{\footnotesize /bin/bash -c "\$(curl -fsSL https://raw.githubusercontent.com/Homebrew/install/HEAD/install.sh)"}
\end{term}
At the end it prints two or three ``Next steps'' lines starting with \texttt{echo} and \texttt{eval} --- run them, then close and reopen Terminal.

\section{Install Node.js, Git and TeX}
\begin{term}
brew install node git\\
brew install --cask mactex-no-gui
\end{term}
\texttt{mactex-no-gui} is the full TeX Live (about 5\,GB) and compiles every template in Numenor. It asks for your Mac password.
When it finishes, \textbf{close and reopen Terminal} so the new \texttt{latexmk} and \texttt{xelatex} commands are found.

\begin{note}[Short on disk space?]
Use the small BasicTeX (about 100\,MB) instead, then add the few tools Numenor needs:
\par\smallskip
\begin{tcolorbox}[colback=navy,colframe=navy,coltext=white,arc=2pt,fontupper=\ttfamily\small,left=6pt,top=3pt,bottom=3pt]
brew install --cask basictex\\
sudo tlmgr update --self\\
sudo tlmgr install latexmk collection-fontsrecommended
\end{tcolorbox}
If a template later reports a missing \texttt{.sty} file, install that package with \texttt{sudo tlmgr install <name>}.
\end{note}

Check that everything is in place (each should print a version):
\begin{term}
node -v \&\& git --version \&\& latexmk -v | head -1
\end{term}
Node must be version 18 or newer.

\section{Download and build Numenor}
\begin{term}
cd \textasciitilde\\
git clone https://github.com/farhadsh1992/numenor.git\\
cd numenor\\
npm install\\
npm run build
\end{term}
This puts Numenor in a \texttt{numenor} folder in your home folder. \texttt{npm install} takes a minute or two the first time.

\section{Start it}
\begin{term}
cd \textasciitilde/numenor \&\& NUMENOR\_HOST=127.0.0.1 npm start
\end{term}
When Terminal shows \textit{Numenor started at http://localhost:8787}, open
\href{http://localhost:8787}{\texttt{http://localhost:8787}} in Safari, Chrome or any browser.
Keep the Terminal window open while you work; press \texttt{Ctrl + C} in it to stop Numenor.

\texttt{NUMENOR\_HOST=127.0.0.1} keeps the editor reachable from your Mac only. Leave it out if you want other
computers on your network to open it. Sharing with people over the internet does not need it removed ---
use the \textbf{Share} button instead (see below).

\begin{note}[Start it with one word]
Add a shortcut once:
\par\smallskip
\begin{tcolorbox}[colback=navy,colframe=navy,coltext=white,arc=2pt,fontupper=\ttfamily\small,left=6pt,top=3pt,bottom=3pt]
{\scriptsize echo "alias numenor='cd \textasciitilde/numenor \&\& NUMENOR\_HOST=127.0.0.1 npm start'" >> \textasciitilde/.zshrc}
\end{tcolorbox}
Open a new Terminal window and type \texttt{numenor} from then on.
\end{note}

\section{First steps in the app}
\begin{itemize}[leftmargin=1.2em,itemsep=0.2em]
\item \textbf{New project} on the Projects page: pick a template (ACL, CVPR, NeurIPS, ICML \dots), or import a \texttt{.zip} or an arXiv source.
\item Opening a project compiles it and shows the PDF. \textbf{Cmd/Ctrl + Shift + J} jumps from the source to the PDF; double-click the PDF to jump back.
\item \textbf{AI help} (optional): click the \textbf{Numenor} name at the top left of the editor to open Settings and enter
an OpenAI-compatible endpoint, API key and model.
\item Projects are saved in \texttt{\textasciitilde/numenor/data}. The Projects page can add other folders, for example one inside Google Drive or iCloud.
\end{itemize}

\section{Add a shared project to your Numenor}
If you installed Numenor, keep projects other people share with you on your own Projects page:
\begin{enumerate}[leftmargin=1.4em,itemsep=0.2em]
\item On the Projects page, click \textbf{Add shared link} (in the left sidebar, under \emph{Import project}).
\item Paste the invite link you were sent --- for example the EDGE-AI-Overleaf link above --- and click \textbf{Add shared project}.
\item The project now stays under \textbf{Shared with me}, which shows whether its owner's computer is reachable.
Click \textbf{Open} to work on it. If the owner sends a new link later, add it the same way; it replaces the old one.
\end{enumerate}

\section{Optional extras}
\begin{itemize}[leftmargin=1.2em,itemsep=0.2em]
\item \textbf{Share a project with people}: install Cloudflare's tunnel tool once, then use \textbf{Share} in the editor.
\begin{term}
brew install cloudflared
\end{term}
\item \textbf{Spelling and grammar} uses macOS's own checker and is only available when Numenor runs inside Maziar.
\end{itemize}

\section{Update to the newest version}
\begin{term}
cd \textasciitilde/numenor \&\& git pull \&\& npm install \&\& npm run build
\end{term}
Then start it again. Your projects in \texttt{data} are not touched.

\section{If something goes wrong}
\begin{description}[leftmargin=1.2em,style=nextline,itemsep=0.3em]
\item[\texttt{command not found: brew / node / latexmk}] Close and reopen Terminal. For Homebrew, run the
\texttt{eval} line it printed at the end of its install.
\item[\textit{Address already in use} / port 8787 busy] Numenor is already running somewhere, or another app uses the port.
Start on another port: \texttt{PORT=8790 npm start}, then open \texttt{http://localhost:8790}.
\item[Compile fails with a missing \texttt{.sty} file] With BasicTeX: \texttt{sudo tlmgr install <package>}. With full MacTeX this should not happen.
\item[The page is blank after an update] Run \texttt{npm run build} again, then reload the browser page.
\end{description}

\vfill
{\small\color{gray} Numenor is based on OpenPrism by OpenDCAI (MIT License). Source and README:
\url{https://github.com/farhadsh1992/numenor}}
\end{document}
